\section{因果前沿与整个候选图的接口}\label{sec:frontier}
一个已处理事件可能仍与许多未来事件相连。若这些连接都指向同一个旧写，保留这个写就足够；若许多旧事件分别有未处理连接，接口才会增长。本节用已经认证的共同因果关系区分这两种情况，并同时处理尚未选择的来源和写序。

\subsection{可靠关系与可检查的来源过滤}
证书 $\Gamma_I$ 包含每个读的候选域 $D(r)$、共同因果关系 $H$、逐地址写序下界 $L$、推导记录，以及原始变量的归属信息。对每个合法完成 $J$，它满足
\begin{equation}\label{eq:reliable}
 \rf_J^{-1}(r)\in D(r),\qquad H\subseteq\hb_J,\qquad L\subseteq\co_J.
\end{equation}
这里 $H$ 和 $L$ 都是严格传递关系，$H$ 包含 $\SB$ 和初始化先行关系，$L$ 只含同地址写对，并包含 $H$ 的同地址写限制及初始化最小关系。$D$ 是所有合法来源的上近似；其中每个元素未必最终有支持。

我们要求证书对下列必要条件达到饱和。若 $H(r,w)$，从 $D(r)$ 删除 $w$；否则选择该来源会形成 $\hb$ 环。若存在同地址写 $z$，使 $L(w,z)$ 且 $H(z,r)$，也删除 $w$；否则 $r\mathrel{\fr}z\mathrel{\hb}r$ 违反一致性。下文分别称这两条为 R1 和 R2。若 $D(r)=\{w\}$，则把 $(w,r)$ 加入 $H$，称为 R3。共同前序规则为
\begin{equation}\label{eq:rall}
 \frac{D(r)\ne\varnothing\quad
       \forall w\in D(r).\ (u=w\ \lor\ H(u,w))}{H(u,r)}\quad(R_{\forall}).
\end{equation}
该规则只引用此前得到的事实。任取合法完成，它为 $r$ 选择的来源 $w$ 仍在 $D(r)$；于是 $u$ 先于该写或等于该写，再经实际 $\rf(w,r)$ 到达 $r$。这证明结论对所有完成成立，且没有把任何候选选成必选来源。

R1--R3 和式~\eqref{eq:rall} 的可靠性属于关系分析的基本推理~\cite{haas2023}。这里需要它们提供的是下一定理的具体前提：稳定写若已被更晚的稳定写覆盖，其所有未来来源候选必须已经删除。只保存来源域而不实施 R2，不能得到这个结构结论。

\subsection{前沿、稳定写与真实候选边界}
令 $\theta$ 为所有非初始化事件的一个线性排列，保持每个线程的 $\SB$。它决定编译时引入事件的次序。令 $P_t$ 为前 $t$ 个事件，$Q_t$ 为其余事件。两个事件称为\emph{相关}，当且仅当它们访问同一地址且至少一个是写。定义
\begin{equation}\label{eq:frontier}
 U_t=\{e\in P_t:\exists f\in Q_t,\ e,f\text{ 相关且 }\neg H(e,f)\},
 \qquad c=\max_t |U_t|.
\end{equation}
$U_t$ 是未决因果前沿。事件 $e\notin U_t$ 表示其与每个未来相关访问的先后已经确定；即使 $e$ 很早出现，也可能仍是未来读取唯一需要的旧写。

对地址 $x$，令 $S_x(t)$ 为 $P_t\setminus U_t$ 中访问 $x$ 的写，令 $M_x(t)$ 为其中关于 $H$ 的极大元素。记 $M_t=\bigcup_x M_x(t)$，并以 $T_t$ 表示每个线程在 $P_t$ 中的最后一个事件。一个线程内的同地址写按 $H$ 全序排列，因此
\begin{equation}\label{eq:maxwrites}
 |M_x(t)|\le p,\qquad |M_t|\le ap,\qquad |T_t|\le p.
\end{equation}
极大性在这里相对于\emph{稳定写}，不是相对于所有已引入写。一个仍在前沿中的新写不能用来代替更老的稳定写。

构造无向候选事件图 $G_I$。它含全部 $w\in D(r)$ 对应的 RF 边、$L$ 中不可比较的每个同地址写对、$\Cov(L)$ 的边、真实线程邻接边，以及每个初始化写到各非空线程首事件的边。不可比较写对以后选择两个 CO 方向之一；$\Cov(L)$ 为固定 CO 提供稀疏生成边。$H$ 和 $L$ 的稠密闭包用于推理，但不作为候选图的全部边加入。令
\[
 F_t=\{e\in P_t:\exists f\in Q_t,\ \{e,f\}\in E(G_I)\}
\]
为非初始化事件的真实候选边界。它直接决定哪些旧事件还有未处理关联。

\begin{theorem}[因果前沿控制统一候选图]\label{thm:frontier}
设 $I$ 含 $n$ 个事件、$p$ 个线程和 $a$ 个地址。设可靠证书 $\Gamma_I$ 满足上述传递性、包含关系及 R1/R2 饱和条件，$\theta$ 为预先固定的 $\SB$ 相容排列，$c$ 由式~\eqref{eq:frontier} 定义。则每个切口都有
\begin{equation}\label{eq:containment}
 F_t\subseteq U_t\cup M_t\cup T_t.
\end{equation}
从 $G_I$ 可直接构造一个覆盖全部候选二元选择的连续袋计划 $\pi_I$，其最大事件袋大小满足
\begin{equation}\label{eq:mainK}
 K\le c+ap+p+a+1.
\end{equation}
这个图和计划同时适用于 $I$ 的所有合法完成及所有后续原始变量查询。
\end{theorem}
\begin{proof}
固定切口，考虑 $e\in F_t\setminus U_t$ 的一条跨切口边 $\{e,f\}$。对于 RF 边有两种端点角色。若 $e$ 为过去的读而 $f$ 为未来的写，稳定性给出 $H(e,f)$，R1 已排除这个来源，所以该边不可能存在。若 $e=w$ 为过去的写、$f=r$ 为未来的读，并且 $w\notin M_t$，有限严格偏序中存在同地址稳定写 $z$，满足 $H(w,z)$。由证书包含关系有 $L(w,z)$；由 $z$ 的稳定性有 $H(z,r)$。R2 已删除 $w\to r$，同样矛盾。因此任何仍与未来读相连的稳定写都在 $M_t$ 中。

再考虑两个写端点 $w=e$、$v=f$。稳定性给出 $H(w,v)$，从而 $L(w,v)$；它们不可能是未决 CO 对。如果 $w\notin M_t$，仍取稳定写 $z$ 使 $H(w,z)$。因为 $v$ 在未来，有 $H(z,v)$，所以 $L(w,z)$ 和 $L(z,v)$ 同时成立。$z$ 严格位于二者之间，排除了 $(w,v)\in\Cov(L)$。故固定覆盖边的旧稳定端点也必须属于 $M_t$。这一步使理论允许跨很长历史的覆盖边：它保留少量极大写，而不要求每条边都短。

真实 $\SB$ 邻接边跨切口时，过去端点必为该线程在 $P_t$ 中的尾事件。因为 $\theta$ 保持线程顺序，不能已有另一个同线程事件越过这个未引入的直接后继。故该端点在 $T_t$ 中。初始化端点单独计入 $a$，不属于 $F_t$ 的计数。所有候选边类型已经覆盖，式~\eqref{eq:containment} 成立。

构造时先引入初始化写并连续保留。在引入 $\theta(t+1)$ 前保留真实边界 $F_t$；引入新事件后，处理以它为较晚端点的所有二元组，随后退休已经没有未来候选关联的事件。当前袋包含在
\[
 \{\text{初始化写}\}\cup F_t\cup\{\theta(t+1)\}
\]
中。结合式~\eqref{eq:maxwrites} 即得式~\eqref{eq:mainK}。每条边在其较晚端点引入时得到处理；每个事件从首次引入至最后关联连续存在。即使某个事件在不同切口的 $U_t$、$M_t$ 分类改变，也不会退休后重新引入，因为实际退休依据是 $F_t$。读的来源计数在其所有 RF 组处理完后检查。全部选择组、最后关联位置和计数检查位置只依赖证书与 $\theta$，不依赖选择了哪个完成或查询。\qedhere
\end{proof}

定理区分了两种压缩。关系推理删去永远不可能的候选；稳定写论证进一步说明，即便某些旧候选一直存在，整个过去仍只需少量代表与未来相连。变量身份随后保留在编译路径中，而事件关系可在接口外被摘要化。第~\ref{sec:compilation}节证明两者兼容。

\subsection{二元选择、阶段分隔与量词边界}
二元 RF 组只处理一个候选 $w\to r$ 的取舍，并在 $r$ 活跃期间累计已选来源数。若把一个读及其全部候选作为一次原子选择，它们必须同时出现在袋中，这超出了定理的二元边覆盖前提。例如一线程依次写 $w_1,\ldots,w_h$，另一线程只有读 $r$，所有写值都等于读值。在先引入 $r$ 的布局上，$c=1,p=2,a=1$；全部来源仍可行，一个原子 RF 组却有 $h+1$ 个非初始化端点。若另给来源域上界 $d=\max_r|D(r)|$，通过延长候选写生存期可得保守原子组界
\[
 K_{\rm atomic}\le c(d+1)+ap+p+a+1.
\]
附录~\ref{app:atomic} 给出相应调度证明。本文的主构造始终采用二元组。

许多协议更容易先证明按阶段的共同因果分隔。若阶段布局保持 $\SB$，每阶段最多 $b$ 个事件，且相隔至少 $h$ 个阶段的事件均有 $H(e,f)$，则切口的 $U_t$ 只可能来自最近 $h$ 个阶段，故 $c\le hb$。定理立即给出 $K\le hb+ap+p+a+1$。已有完整交替发布网络的来源块锁定证明可以产生 $h=2d+1$，其中 $d$ 为通信网络直径。由此，其接口界不再需要“每隔 $G$ 阶段必有一次写”的写覆盖假设：一个地址只在开始和很晚以后写入，中间反复读取，也只需保留稳定的极大写。该推论保留了原网络的完整交替观测前提；第~\ref{sec:protocols}节通过真实反馈给出适用于有限前缀的另一种分隔来源。

不能用“每个完成都存在小分解”替代定理中的同一候选结构。考虑一线程对同一地址写 $h$ 次相同值，另一线程连续读 $h$ 次该值。任何合法完成的来源编号按 RR 一致性非降。按写编号依次引入该写及从它读取的读，可用至多三个非初始化端点覆盖该完成的线程邻接与已选 RF 图。可是每对写、读都在某个 SC 完成中相连：把全部读放在指定写之后、下一写之前即可。因此全候选 RF 图包含 $K_{h,h}$，其任意路径分解的袋至少为 $h+1$。这展示的是量词顺序差异：逐完成选择布局不能控制候选并图。本文的协议论证必须在选择来源之前构造 $\Gamma_I,G_I,\pi_I$。
